\paragraph{Evasion.} Every learned descriptor could be evaded by a white-box attacker with an imperceptible perturbation (Fig.~\ref{fig:fp-security}a): at a budget of 2 grey levels (PSNR \val{fp:sec/whitebox/DINOv2-S/2/psnr}\,dB) projected gradient descent pushed \val{fp:sec/whitebox/DINOv2-S/2/pct7}\% of the registered images below the $10^{-7}$ threshold for DINOv2-S, and equally for SSCD, ISC21, OpenCLIP and the keyed projection, whose soft pre-sign values carry the gradient. The DCT hashes resisted longer but not for long: through differentiable surrogates the attack evaded PDQ for \val{fp:sec/surrogate/PDQ/2/pct7}\%, \val{fp:sec/surrogate/PDQ/4/pct7}\% and \val{fp:sec/surrogate/PDQ/8/pct7}\% of images at 2, 4 and 8 grey levels, pHash for \val{fp:sec/surrogate/pHash-64/2/pct7}\% already at 2, and the ISCC Image-Code for \val{fp:sec/surrogate/ISCC-Image-64/2/pct7}\% at 2. A black-box attacker who crafts the perturbation on an ensemble of DINOv2-S, SSCD and OpenCLIP (Fig.~\ref{fig:fp-security}b, budget 8) evades every DINOv2 variant and SSCD-R50 but transfers poorly to ISC21 (\val{fp:sec/transfer/8/ISC21-1st/pct}\%) and not at all to PDQ and pHash-256; the ISCC codes are evaded for about half of the images at every budget (\val{fp:sec/transfer/2/ISCC-Image-64/pct}\% at 2 grey levels), again through the border-trimming step, which small perturbations of a white background switch on or off.

\paragraph{Forgery.} Binding an unrelated image to a registered asset is harder but, for the learned copy detectors, practical (Fig.~\ref{fig:fp-security}c): at 4 grey levels the attack made \val{fp:sec/collision/ISC21-1st/4/pct}\% of source images the top match of their target in the full registry above threshold for ISC21 and \val{fp:sec/collision/SSCD-mixup/4/pct}\% for SSCD-R50, rising to \val{fp:sec/collision/ISC21-1st/8/pct}\% and \val{fp:sec/collision/SSCD-mixup/8/pct}\% at 8. DINOv2-S (\val{fp:sec/collision/DINOv2-S/8/pct}\% at 8) and OpenCLIP (\val{fp:sec/collision/OpenCLIP-B32/8/pct}\%) were harder to steer, and the hashes nearly impossible within these budgets (PDQ \val{fp:sec/collision/PDQ/16/pct}\% and ISCC \val{fp:sec/collision/ISCC-Image-64/16/pct}\% at 16 grey levels, where the PSNR is already \val{fp:sec/collision/PDQ/16/psnr}\,dB). The copy-detection descriptors, which generalise best across benign transformations, are thus also the easiest to forge: their similarity is smooth in the pixels, which is what both robustness and gradient steering exploit.

\paragraph{Inversion.} A soft-binding value is published inside the manifest, so we asked how much of the image it gives away. A decoder trained on the descriptors of the \val{fp:sec/inv/ntrain} distractors reconstructed a $64\times64$ thumbnail of each of \val{fp:sec/inv/ntest} registered images from its descriptor alone (Fig.~\ref{fig:fp-inversion}). Reconstructions from the float descriptors were recognisable in colour, layout and object class (LPIPS \val{fp:sec/inv/DINOv2-S/lpips} for DINOv2-S against \val{fp:sec/inv/mean-image/lpips} for the mean image), and using them as queries retrieved the true image among all \val{fp:sec/inv/ntest} in \val{fp:sec/inv/OpenCLIP-B32/reid/pct}\% (OpenCLIP), \val{fp:sec/inv/DINOv2-S/reid/pct}\% (DINOv2-S) and \val{fp:sec/inv/SSCD-mixup/reid/pct}\% (SSCD) of cases, 40 to 95 times chance. The binary hashes leaked less but still identified images above chance (PDQ \val{fp:sec/inv/PDQ/reid/pct}\%). The keyed projection behaved as intended: an attacker who knows the key still identifies images at about 40 times chance (\val{fp:sec/inv/DINOv2-S-LSH256/reid/pct}\%, half the rate from the float DINOv2-S descriptor), but a decoder trained with any other key reconstructs only the mean image (LPIPS \val{fp:sec/inv/DINOv2-S-LSH256-wrongkey/lpips}) and identifies images at chance (\val{fp:sec/inv/DINOv2-S-LSH256-wrongkey/reid/pct}\%).

\begin{figure*}[!t]
\centering
\includegraphics[width=\textwidth]{fp_fig10_inversion}
\caption{Inverting stored fingerprints. (a) Registered images (top) and thumbnails reconstructed from their descriptors by a decoder trained on distractor descriptors; bottom rows: the keyed 256-bit projection decoded by an attacker with the right and with a wrong key. (b) Share of reconstructions that retrieve their own original among the \val{fp:sec/inv/ntest} registered images (95\% Clopper--Pearson intervals); dashed: chance.}
\label{fig:fp-inversion}
\end{figure*}


\begin{figure*}[!t]
\centering
\includegraphics[width=\textwidth]{fp_fig9_security}
\caption{Adversarial robustness on \val{fp:sec/n} registered ABO images per setting. (a) Evasion by white-box PGD (learned descriptors) and surrogate-gradient PGD (DCT hashes) against the $L_\infty$ budget, judged by the real extractor at pair-level FPR $10^{-7}$. (b) Evasion by perturbations crafted on a DINOv2-S, SSCD-R50 and OpenCLIP ensemble at a budget of 8 grey levels, judged at the pair-level $10^{-6}$ threshold the attack ran against; n/a: the method has no reachable operating point. (c) Targeted forgery: share of unrelated images that become their target's top match above the $10^{-6}$ threshold.}
\label{fig:fp-security}
\end{figure*}

\paragraph{Audio and video.} The same pattern holds in the other modalities. Audio attacks were judged, like the image attacks, at the pair-level $10^{-6}$ threshold of each method, taken from the first audio retrieval run; those thresholds differ from the re-run reported in Section~\ref{sec:fp-res-audio} by at most \val{fp:sec/audio/thdelta}, and re-scoring every CLAP attack at the re-run thresholds changes \val{fp:sec/audio/rescore_changed} outcomes. Video attacks were judged at a different threshold, described below. For audio, projected gradient descent on the CLAP embedding of \val{fp:sec/audio/n} registered tracks evaded CLAP for every track at every budget we tried, down to a target signal-to-noise ratio of 40\,dB (achieved median \val{fp:sec/audio/snr/40}\,dB). The perturbation did not transfer to the recording-level fingerprints: Chromaprint matched every attacked track at all budgets, and the ISCC Audio-Code and audfprint lost \val{fp:sec/audio/pgd/20/ISCC-Audio-64/pct}\% and \val{fp:sec/audio/pgd/20/audfprint/pct}\% only at the loudest budget (20\,dB target), where the added signal acts as noise rather than as a targeted attack (\val{fp:sec/audio/pgd/40/ISCC-Audio-64/pct}\% and \val{fp:sec/audio/pgd/40/audfprint/pct}\% at 40\,dB). An attacker does not need gradients against those systems, however: a search over pitch and tempo changes found the smallest change that broke each match, and the median was \val{fp:sec/audio/mag/Chromaprint/pitch} semitones or a tempo factor of \val{fp:sec/audio/mag/Chromaprint/tempo} for Chromaprint and \val{fp:sec/audio/mag/audfprint/pitch} semitones or \val{fp:sec/audio/mag/audfprint/tempo} for audfprint and the ISCC Audio-Code, where 0.25 semitones and 1.02 are the smallest steps of the grid. Only CLAP needed audible changes (\val{fp:sec/audio/mag/CLAP/pitch} semitones, tempo \val{fp:sec/audio/mag/CLAP/tempo}). Every audio fingerprint in this study is therefore evaded by a small change, either by gradients (CLAP) or by a pitch shift (all others).

For video, frame-wise projected gradient descent against DINOv2-S and SSCD on \val{fp:sec/video/n} registered clips evaded both frame-level neural descriptors, in sequence and mean-pooled form, for every clip at 4 grey levels, and transferred to none of the hashes: vPDQ, TMK+PDQF and both ISCC Video-Codes matched every attacked clip at 4 and 8 grey levels (judged at each method's query-level 1\% threshold). The video calibration set does not resolve a pair-level $10^{-6}$ target, so every video attack used that threshold, which let through \val{fp:video/fprq/min}\% to \val{fp:video/fprq/max}\% of held-out negative queries (Section~\ref{sec:fp-res-video}). The video results therefore do not establish attack performance at a pair-level false-match rate. Each hash's pair-level $10^{-4}$ threshold is lower than its query-level threshold, so the hashes' matches hold at that target too; the per-clip scores of the neural descriptors were not kept, so their evasion was not re-scored at it. Simple edits again did what gradients could not. A centred crop keeping 90\% of the width and height, the first step of the grid, broke the median match for vPDQ, TMK+PDQF, both ISCC Video-Codes and the sequence-matched DINOv2-S; only SSCD and mean-pooled DINOv2-S held at 90\% and broke at a median of \val{fp:sec/video/magnitude/SSCD-mixup-seq/crop_keep_median}. Playback speed-up was the reverse: TMK+PDQF and the 256-bit ISCC Video-Code still matched all \val{fp:sec/video/magnitude/TMK+PDQF/speed_never} clips at twice the speed (the 64-bit code \val{fp:sec/video/magnitude/ISCC-Video-64/speed_never}), because dropping frames leaves the remaining frames' hashes intact, while the sequence-matched neural descriptors broke at a median factor of \val{fp:sec/video/magnitude/DINOv2-S-seq/speed_median}. No video fingerprint survived both a small crop and a white-box attack.

